\chapter{\IfLanguageName{dutch}{Inleiding}{Introduction}}%
\label{ch:inleiding}
Voertuigen bevatten steeds meer elektronische systemen die tijdens het rijden
grote hoeveelheden gegevens verzamelen. Via de OBD-II-interface kan een deel
van deze voertuigdata worden uitgelezen, waaronder gegevens over onder andere
het motortoerental, de voertuigsnelheid, temperaturen en de belasting van de
motor. Deze gegevens kunnen worden gebruikt voor diagnose, analyse en het
realtime opvolgen van de toestand van een voertuig.

Het uitlezen van deze gegevens vormt echter slechts één onderdeel van een
monitoringsysteem. Om voertuigdata op een bruikbare manier aan een gebruiker
te presenteren, moeten de gegevens worden uitgelezen, verwerkt, doorgestuurd
naar een gebruikersinterface en daar op een overzichtelijke manier worden
weergegeven. Hierbij speelt ook de gekozen communicatiemethode tussen backend
en frontend een belangrijke rol, zeker wanneer gegevens frequent moeten worden
bijgewerkt.

Binnen deze bachelorproef wordt daarom een webgebaseerde applicatie ontwikkeld
voor het realtime monitoren van OBD-II-voertuigdata. Daarnaast wordt onderzocht
hoe REST API en WebSockets zich binnen deze toepassing tot elkaar verhouden.
Door beide communicatiemethoden onder dezelfde omstandigheden te testen, wordt
nagegaan welke methode het meest geschikt is voor het doorsturen van realtime
voertuigdata.

\section{Probleemstelling}
Moderne voertuigen verzamelen via verschillende elektronische regeleenheden
een grote hoeveelheid technische informatie. Een deel van deze gegevens is via
de OBD-II-interface toegankelijk voor externe toepassingen. Hoewel OBD-II een
gestandaardiseerde interface biedt, verschilt de daadwerkelijk beschikbare
informatie tussen voertuigen. Niet ieder voertuig ondersteunt namelijk dezelfde
parameters.

Daarnaast moet de uitgelezen voertuigdata op een efficiënte manier van de
backend naar de frontend van een monitoringapplicatie worden doorgestuurd.
Hiervoor bestaan verschillende communicatiemethoden. REST API's worden veel
gebruikt binnen webapplicaties en werken volgens een request-responseprincipe.
Voor realtime toepassingen kan hierbij polling worden gebruikt, waarbij de
frontend op vaste tijdstippen nieuwe gegevens opvraagt. WebSockets gebruiken
daarentegen een blijvende verbinding waarlangs opeenvolgende gegevens kunnen
worden verstuurd.

Bij frequent veranderende voertuigdata is het daarom relevant om na te gaan
welke invloed deze verschillende communicatieprincipes hebben op de overdracht
van de gegevens. Vooral latency en netwerkverbruik zijn hierbij van belang.
Ook de gekozen updatefrequentie kan invloed hebben op deze eigenschappen.

Een bijkomend praktisch probleem is dat tijdens de ontwikkeling niet voortdurend
een fysiek voertuig met een aangesloten OBD-II-adapter beschikbaar is. Daarom
is een simulatieomgeving nuttig waarmee de werking van de applicatie ook zonder
fysieke voertuigverbinding kan worden ontwikkeld en getest.

Deze bachelorproef richt zich daarom op de ontwikkeling van een volledige
monitoringapplicatie, vanaf het uitlezen van OBD-II-data tot de realtime
weergave ervan in een webinterface. Daarbij wordt specifiek onderzocht hoe
REST API en WebSockets presteren bij verschillende updatefrequenties.

\section{Onderzoeksvraag}
De centrale onderzoeksvraag van deze bachelorproef luidt als volgt:

\begin{quote}
\textit{Hoe kan OBD-II-voertuigdata realtime worden uitgelezen, verwerkt en
gevisualiseerd in een webapplicatie, en welke van de communicatiemethoden REST
API en WebSockets is binnen deze toepassing het meest geschikt voor het
doorsturen van deze data?}
\end{quote}

Om deze hoofdvraag te beantwoorden, worden de volgende deelvragen onderzocht:

\begin{itemize}
    \item Welke voertuigparameters kunnen via OBD-II worden uitgelezen en welke
    verschillen bestaan hierbij tussen de onderzochte voertuigen?

    \item Hoe kan een webapplicatie worden opgebouwd om OBD-II-voertuigdata
    realtime te verwerken en te visualiseren?

    \item Hoe kunnen voertuiggegevens worden gesimuleerd zodat de applicatie ook
    zonder een fysieke voertuigverbinding kan worden ontwikkeld en getest?

    \item Welke verschillen bestaan tussen REST API en WebSockets op het vlak
    van latency en netwerkverbruik bij update-intervallen van 100, 250 en
    1000 milliseconden?

    \item Welke communicatiemethode en updatefrequentie zijn binnen de
    ontwikkelde toepassing het meest geschikt voor realtime
    voertuigmonitoring?
\end{itemize}


\section{Onderzoeksdoelstelling}
Het doel van deze bachelorproef is het ontwikkelen en evalueren van een
webgebaseerde applicatie voor het realtime monitoren van voertuigdata via
OBD-II.

Eerst wordt onderzocht welke voertuigparameters via OBD-II beschikbaar zijn.
Hiervoor worden meerdere voertuigen getest en wordt nagegaan welke parameters
gemeenschappelijk beschikbaar zijn en welke verschillen tussen voertuigen
optreden.

Vervolgens wordt een monitoringapplicatie ontwikkeld die bestaat uit een
backend voor het uitlezen en verwerken van voertuigdata en een frontend voor
het realtime weergeven ervan. Daarnaast wordt een simulatieomgeving ontwikkeld
waarmee voertuigdata kan worden nagebootst wanneer geen fysiek voertuig
beschikbaar is.

Binnen de ontwikkelde applicatie worden zowel REST API als WebSockets
geïmplementeerd. Beide communicatiemethoden worden vervolgens getest met
update-intervallen van 100, 250 en 1000 milliseconden. Hierbij worden de
latency en het netwerkverbruik gemeten en met elkaar vergeleken.

Het eindresultaat moet aantonen in welke mate de ontwikkelde architectuur
geschikt is voor realtime OBD-II-monitoring en welke van de onderzochte
communicatiemethoden binnen deze toepassing de meest geschikte oplossing vormt.


\section{Opzet van deze bachelorproef}
De rest van deze bachelorproef is als volgt opgebouwd.

In Hoofdstuk~\ref{ch:stand-van-zaken} wordt de bestaande kennis rond OBD-II,
realtime dataverwerking, communicatietechnologieën, simulatie en visualisatie
besproken. Dit hoofdstuk vormt de theoretische basis voor de verdere
ontwikkeling en vergelijking.

In Hoofdstuk~\ref{ch:methodologie} wordt de gevolgde onderzoeksaanpak
toegelicht. De verschillende fasen van het onderzoek worden beschreven, van
het verzamelen van voertuigdata tot de ontwikkeling en evaluatie van de
monitoringapplicatie.

In Hoofdstuk~\ref{ch:onderzoek-en-verzameling-van-voertuigdata} wordt
beschreven hoe de verbinding met verschillende voertuigen tot stand wordt
gebracht en welke OBD-II-parameters beschikbaar zijn. De resultaten worden
vergeleken om te bepalen welke gegevens verder binnen de applicatie kunnen
worden gebruikt.

In Hoofdstuk~\ref{ch:ontwikkeling-van-een-obd-ii-simulatieomgeving} wordt de
ontwikkeling van de simulatieomgeving besproken. Deze omgeving maakt het
mogelijk om de applicatie ook zonder een aangesloten voertuig verder te
ontwikkelen en te testen.

In Hoofdstuk~\ref{ch:ontwikkeling-van-een-realtime-monitoringdashboard} wordt
de ontwikkeling van de monitoringapplicatie beschreven. Hierbij komen de
architectuur, backend, frontend en datastroom tussen de verschillende
componenten aan bod.

In Hoofdstuk~\ref{ch:vergelijking-van-rest-api-en-websockets} worden REST API
en WebSockets met elkaar vergeleken. Beide communicatiemethoden worden getest
bij verschillende update-intervallen en geëvalueerd op basis van latency en
netwerkverbruik.

Vervolgens wordt het volledige systeem getest en geëvalueerd. Hierbij wordt
nagegaan in welke mate de ontwikkelde applicatie voldoet aan de vooropgestelde
doelstellingen en worden de belangrijkste beperkingen van het onderzoek
besproken.

Tot slot wordt in Hoofdstuk~\ref{ch:conclusie} een antwoord geformuleerd op
de centrale onderzoeksvraag en de bijhorende deelvragen. De belangrijkste
bevindingen worden samengevat en er worden mogelijkheden voor toekomstig
onderzoek besproken.